% mode: ini
% Aimed at after_math's display code, 1377-1383: e, q, the position d and the too-close rule; the glue
% chosen by d+s against \predisplaysize; \eqno and \leqno, each with the number on its own line or
% in the line; with glue first in the formula, which sets d to 0; \vadjust, \vadjust pre and \insert
% coming out at the right place; \predisplaypenalty and \postdisplaypenalty; the code g2=0.
\input prelude
\def\e#1{\vrule width#1pt height1pt}
\def\sh{\showbox1 }
\def\d#1{\setbox1\vbox{\noindent\predisplaysize=\ps $$#1$$}\sh}
\dimendef\ps=200 \ps=10pt
\hsize=100pt
\abovedisplayskip=11pt \belowdisplayskip=12pt \abovedisplayshortskip=13pt \belowdisplayshortskip=14pt
\predisplaypenalty=21 \postdisplaypenalty=22
% No number: long and short skips by d+s<=\predisplaysize.
\ps=-16383.99999pt \d{\e{10}}
\ps=44pt \d{\e{10}}
\ps=45pt \d{\e{10}}
\ps=46pt \d{\e{10}}
\ps=44pt \d{\displayindent=5pt \e{10}}
\ps=50pt \d{\displayindent=5pt \e{10}}
\ps=-50pt \d{\displayindent=-5pt \displaywidth=50pt \e{10}}
\ps=-16383.99999pt
% The odd half: z-w odd, then negative, half rounds away by half's own rule.
\setbox1\vbox{\noindent$$\displaywidth=11sp \e{0}\kern 4sp$$}\sh
\setbox1\vbox{\noindent$$\displaywidth=11sp \e{0}\kern 5sp$$}\sh
\setbox1\vbox{\noindent$$\displaywidth=11sp \kern 12sp$$}\sh
\setbox1\vbox{\noindent$$\displaywidth=11sp \kern 13sp$$}\sh
% \eqno and \leqno, plenty of room.
\d{\e{10}\eqno\e{5}}
\d{\e{10}\leqno\e{5}}
\d{\eqno\e{5}}
\d{\leqno\e{5}}
\d{\e{10}\eqno}
\d{\e{10}\leqno}
\d{\eqno}
\d{\leqno}
% \mathsurround does not affect it; math glue in the number.
\mathsurround=2pt \d{\e{10}\eqno\e{5}} \mathsurround=0pt
\d{\e{10}\eqno\hskip 3pt plus 1fil\e{5}}
\d{\e{10}\eqno\mskip 18mu\e{5}}
% Too close: d<2*e, so d is half(z-w-e), or 0 if the display starts with glue.
\d{\e{30}\eqno\e{20}}
\d{\e{30}\leqno\e{20}}
\d{\hskip 3pt\e{27}\eqno\e{20}}
\d{\hskip 3pt\e{27}\leqno\e{20}}
\d{\e{1}\hskip 3pt\e{27}\eqno\e{20}}
\d{\mskip 18mu\e{27}\eqno\e{20}}
\d{\mathsurround=3pt\hskip 3pt\e{27}\eqno\e{20}}
\d{\mathchar"141 \eqno\e{49}}
\d{\kern 3pt\e{27}\eqno\e{20}}
\d{\penalty 3\hskip 3pt\e{27}\eqno\e{20}}
\d{\hbox{}\e{27}\eqno\e{20}}
\d{\e{0}\eqno\e{20}}
% e at 0: number of width 0, and of width below it; it is not treated as a number.
\d{\e{10}\eqno\kern 0pt}
\d{\e{10}\leqno\kern 0pt}
\d{\e{10}\eqno\kern-3pt}
\d{\e{10}\leqno\kern-3pt}
\d{\e{60}\eqno\kern-3pt}
\d{\e{98}\eqno\kern-3pt}
% Adjust material: from the formula.  A vadjust in the number has nowhere to go: it is in a's box.
\skip100=0pt \count100=1000 \dimen100=100pt
\d{\e{10}\vadjust{\kern 1pt}}
\d{\e{10}\vadjust pre{\kern 2pt}}
\d{\e{10}\vadjust{\kern 1pt}\vadjust pre{\kern 2pt}\eqno\e{5}}
\d{\e{10}\vadjust{\kern 1pt}\vadjust pre{\kern 2pt}\leqno\e{5}}
\d{\e{98}\vadjust{\kern 1pt}\vadjust pre{\kern 2pt}\eqno\e{5}}
\d{\e{98}\vadjust{\kern 1pt}\vadjust pre{\kern 2pt}\leqno\e{5}}
\d{\e{10}\eqno\e{5}\vadjust{\kern 1pt}\vadjust pre{\kern 2pt}}
\d{\e{10}\insert100{\kern 3pt}\mark{a}}
\d{\e{98}\insert100{\kern 3pt}\mark{a}\eqno\e{5}}
\d{\e{10}\mark{a}\eqno\mark{b}\e{5}}
\d{\e{10}\vadjust{\kern 1pt}\vadjust{\kern 3pt}\vadjust pre{\kern 2pt}\vadjust pre{\kern 4pt}}
% Display in a display's box, and vcenter; entries.  \halign in an entry of the formula.
\d{\vcenter{\hbox{$$\e{5}$$}}}
\d{\hbox{\e{5}$\e{3}$}}
% Penalties and skip codes in their odd values.
\predisplaypenalty=10000 \postdisplaypenalty=-10000 \d{\e{10}}
\predisplaypenalty=21 \postdisplaypenalty=22
\abovedisplayskip=0pt \belowdisplayskip=0pt \abovedisplayshortskip=0pt \belowdisplayshortskip=0pt \d{\e{10}}
\abovedisplayskip=11pt plus 1pt minus 1pt \d{\e{10}\eqno\e{5}}
% \parshape and \hangindent: s, in d+s.
\ps=50pt
\setbox1\vbox{\hangindent=5pt\hangafter=0 \noindent\predisplaysize=\ps $$\e{10}$$}\sh
\setbox1\vbox{\hangindent=5pt\hangafter=0 \noindent\predisplaysize=\ps $$\e{10}\eqno\e{5}$$}\sh
\ps=-16383.99999pt
% \eqno with an \everymath.
\everymath={\e{1}}
\d{\e{10}\eqno\e{5}}
\everymath={}
\everydisplay={\e{1}}
\d{\e{10}\eqno\e{5}}
\everydisplay={}
\end
